\documentclass[11pt,a4paper]{article}
\usepackage[utf8]{inputenc}
\usepackage[T1]{fontenc}
\usepackage{amsmath}
\usepackage{amssymb}
\usepackage{graphicx}
\usepackage{booktabs}
\usepackage{hyperref}
\usepackage{geometry}
\geometry{margin=0.8in}

\title{Understanding Context Sampling in TabPFN\\ on Small Tabular Datasets}

\author{
  Mohammed Abdullah\\
  \texttt{mohammed.a@bigbangboom.com}
}

\date{\today}

\begin{document}

\maketitle

\begin{abstract}
TabPFN classifies by in-context learning: it conditions on a set of labeled
training rows (the context, or prototypes) and predicts test labels without
gradient updates. On small tabular datasets, where TabPFN is designed to be
used, the practitioner still chooses how large that context is and which rows
go into it. We ask how much those choices matter. Using repeated random
sub-sampling on six small OpenML datasets, we study three linked questions:
(i) how variable predictions are across different random contexts of the same
size, and whether larger contexts are more stable; (ii) whether accuracy tracks
how well a sampled context preserves the full training distribution; and
(iii) whether more expensive selection (K-Means, farthest-point sampling) buys
anything over uniform random selection once the context is representative. We
report only measured results and treat them as observations. % TODO fill headline
This is a small-scale study (six datasets, CPU), so we describe conclusions as
intuitions rather than established laws.

\textbf{Keywords:} in-context learning, TabPFN, prototype sampling, context
stability, small data
\end{abstract}

\section{Introduction}

An in-context learner like TabPFN takes a set of labeled examples as input and
produces a predictor in a single forward pass. The set it conditions on---the
context---is chosen at inference time: how many rows ($k$) and which rows. On
large datasets the interesting question is how to fit within a compute budget.
On \emph{small} datasets, the regime TabPFN targets, the question is different
and more basic: given that we could pass all the data, how much does the
\emph{choice} of context actually change the model's behavior?

We investigate three connected hypotheses on six small OpenML datasets
(all with fewer than about 1{,}000 training rows):

\begin{description}
\item[H1 (Stability).] Different random contexts of the same size $k$ yield
different predictions; the spread of accuracy across random draws shrinks as $k$
grows.
\item[H2 (Representativeness).] A context that better preserves the full
training distribution yields higher accuracy; this is why uniform random
sampling works well.
\item[H3 (Selection cost).] Once a context is representative, expensive
selection methods add little over random selection.
\end{description}

These are mechanistic questions about \emph{how} TabPFN responds to its context,
not benchmark claims, so they can be studied with a few well-chosen datasets and
many repeated draws rather than a large dataset suite. We also report two
supporting stress tests: duplicating part of the context (redundancy) and
dropping part of it (robustness).

We state limitations plainly. Six datasets is a small sample; runs are CPU-only;
and we do not prove mechanisms, only measure behavior consistent with them.

\section{Setup}

We use TabPFN v3 with \texttt{ignore\_pretraining\_limits=True}. Each dataset is
split 80/20 (stratified, fixed split seed). A context of size $k$ is formed by
selecting $k$ training rows and fitting TabPFN on exactly those rows
(\texttt{clf.fit(X[idx], y[idx])}); the held-out 20\% is the test set throughout.
We report ROC AUC (one-vs-rest macro for the multiclass dataset). For the
stability and representativeness analyses we draw \textbf{20 independent random
contexts} per $k$; for selection-method and stress tests we use three seeds.

\begin{table}[h]
\centering
\caption{Datasets (all small; public OpenML).}
\label{tab:datasets}
\begin{tabular}{lrrr}
\toprule
Dataset & rows & features & classes \\
\midrule
diabetes           & 768  & 8  & 2 \\
credit-g           & 1000 & 20 & 2 \\
vehicle            & 846  & 18 & 4 \\
wdbc               & 569  & 30 & 2 \\
breast-w           & 699  & 9  & 2 \\
blood-transfusion  & 748  & 4  & 2 \\
\bottomrule
\end{tabular}
\end{table}

\textbf{Representativeness metrics.} For a sampled context we measure two simple
distances from the full training set: (i) the total-variation distance between
the context's class distribution and the full class distribution
(\texttt{class\_tv}); and (ii) the mean absolute difference between the
context's standardized per-feature means and the full training means
(\texttt{feat\_shift}). Smaller values mean the context looks more like the
whole training set.

\section{H1: Context Stability}

% TODO: table of AUC mean/std per k per dataset (block=stability_repr).
\begin{table}[h]
\centering
\caption{Accuracy across 20 random contexts per $k$: mean and standard
deviation of ROC AUC. [PLACEHOLDER - fill from exp14 results]}
\label{tab:stability}
\begin{tabular}{lrrrrr}
\toprule
Dataset & $k{=}16$ & $k{=}32$ & $k{=}64$ & $k{=}128$ & $k{=}256$ \\
\midrule
diabetes & -- & -- & -- & -- & -- \\
\bottomrule
\end{tabular}
\end{table}

% TODO figure fig_stability.pdf: std(AUC) vs k, one line per dataset.
\textbf{What we observe.} [PLACEHOLDER: whether std shrinks with k, by how much,
and any dataset that stays noisy.]

\section{H2: Representativeness}

% TODO: correlation between (class_tv, feat_shift) and AUC, pooled per dataset.
\textbf{What we observe.} [PLACEHOLDER: sign and strength of correlation between
distribution distance and AUC; whether the most-representative draws are the
best-performing.]

\section{H3: Selection Cost}

% TODO: table random vs kmeans vs fps at fixed k, AUC + selection time.
\begin{table}[h]
\centering
\caption{Selection methods at fixed $k$: AUC (3 seeds) and selection time.
[PLACEHOLDER]}
\label{tab:selection}
\begin{tabular}{llrr}
\toprule
Dataset & Method & ROC AUC & selection time (s) \\
\midrule
-- & random & -- & -- \\
\bottomrule
\end{tabular}
\end{table}

\textbf{What we observe.} [PLACEHOLDER: whether kmeans/fps beat random, and the
time cost.]

\section{Supporting: Redundancy and Robustness}

% TODO: redundancy (duplicate fraction) and robustness (drop fraction) tables.
\textbf{Redundancy.} [PLACEHOLDER: does duplicating context rows change AUC?]

\textbf{Robustness.} [PLACEHOLDER: how fast does AUC fall as context is dropped?]

\section{Discussion}

[PLACEHOLDER: tie the three hypotheses together at the strength the data
supports.]

\section{Limitations}

\begin{itemize}
\item Six small datasets; results may not extend to larger or higher-dimensional
data.
\item CPU-only; timing comparisons are CPU costs.
\item Representativeness is measured with two coarse distances; richer distances
(e.g.\ full distributional divergences) are future work.
\item We observe behavior consistent with the hypotheses but do not prove the
underlying mechanism.
\end{itemize}

\section*{Reproducibility}

All numbers are read from a single logged run
(\texttt{experiments/results/exp14\_context\_sampling\_*.json}). Datasets are
public OpenML tasks; the split seed and per-block seeds are recorded.

\begin{thebibliography}{9}
\bibitem{hollmann2024tabpfn}
Hollmann, N., Müller, S., Eggensperger, K., \& Hutter, F. (2024).
Accurate predictions on small data with a tabular foundation model.
\textit{Nature}.
\end{thebibliography}

\end{document}
